% ===================== 第 4 章（上） 写作真题 =====================
\section{真题分类汇编：Writing}

\begin{tipbox}[title=本章说明 / Note]
本章为历年写作真题整理，\textbf{题目均以英文原文呈现}。要求以考试当年的题干为准：多数年份给出“首句 / 名言”，需以此开篇续写（150--200 词）。题目信息整理自公开真题资料，按年份与套数排列；\textbf{Sample essays} 为本指南自编范文，供模仿结构使用。
\end{tipbox}

\subsection{Real Exam Prompts (2016--2025)}

\begin{center}
\small
\begin{longtable}{p{2.1cm}p{11.6cm}}
\toprule
\textbf{Date} & \textbf{Prompt / Requirement (abridged)} \\
\midrule
\endfirsthead
\toprule
\textbf{Date} & \textbf{Prompt / Requirement (abridged)} \\
\midrule
\endhead
2016.06 & Imagine what will happen when people spend more and more time in the virtual world instead of interacting in the real world. \\
\midrule
2016.12 & The importance of innovation and measures to be taken to encourage innovation / creation / invention. \\
\midrule
2017.06 & Whether to major in science or humanities at college. \\
\midrule
2017.12 & Comment on the saying ``Respect others, and you will be respected.'' \\
\midrule
2018.06 & The importance of building trust between employers and employees. \\
\midrule
2018.12 & How to balance job responsibilities and personal interests. \\
\midrule
2019.06 & The importance of mutual understanding and respect in interpersonal relationships. \\
\midrule
2019.12 & The importance of having a sense of community responsibility. (Sets also covered family responsibility and collective responsibility.) \\
\midrule
2020.07 & Comment on the saying ``The best preparation for tomorrow is doing your best today.'' \\
\midrule
2020.09 & Comment on the saying ``What is worth doing is worth doing well.'' \\
\midrule
2020.12 & Why students should be encouraged to develop the ability to meet challenges. \\
\midrule
2021.06 & Achievements in higher education / in poverty alleviation / in urbanization. (Three sets.) \\
\midrule
2021.12 & Comment on the phenomenon of star worship / of fake information online / of overprotecting children, and give measures. (Three sets.) \\
\midrule
2022.06 & ``Nowadays more and more people keep learning new skills to adapt to a fast-changing world.'' (Set 1) ``\dots choose to live an environmentally friendly lifestyle.'' (Set 2) ``\dots take delight in helping the needy.'' (Set 3) \\
\midrule
2022.09 & ``Nowadays more and more students are becoming increasingly aware of the importance of developing digital skills.'' (Set 1) ``Today more and more people begin to realize the pleasures and joys of real-world social interaction.'' (Set 2) ``It is now widely accepted that mutual trust and openness is the key to promoting cooperation.'' (Set 3) \\
\midrule
2022.12 & The use of information technology in education. \\
\midrule
2023.06 & Mental well-being / friendly discussion when opinions differ / an important goal of education is to help students learn how to learn. (Three sets.) \\
\midrule
2023.12 & ``With their valuable skills and experience, elderly people can continue to make significant contributions to society.'' (Sets also covered the importance of acquiring basic knowledge.) \\
\midrule
2024.06 & ``There is a growing awareness of the importance of digital literacy and skills in today's world.'' (Sets also covered autonomous learning and the equal importance of social practice and theoretical study.) \\
\midrule
2024.12 & Setting goals and working persistently lead to success / university life is full of possibilities. \\
\midrule
2025.06 & Teachers' influence on students' academic pursuit and personal development / realizing self-value through the Chinese Dream / preparing well for the diverse challenges in university. \\
\midrule
2025.12 & The influence of teachers on students' academic and personal growth / realizing one's value through the Chinese Dream / being well prepared for diverse challenges in university. \\
\bottomrule
\end{longtable}
\end{center}

\subsection{Task Types and Frameworks}

\begin{enumerate}
  \item \textbf{Phenomenon-explanation} (e.g.\ 2016.06, 2018.06, 2019.06): describe the phenomenon, analyze its causes and effects, then give your attitude.
  \item \textbf{Problem-solving} (e.g.\ 2016.12, 2018.12, 2021.12): state the problem, analyze why it matters, propose two or three concrete measures.
  \item \textbf{Opinion-choice} (e.g.\ 2017.06): present both sides briefly, then state and defend your choice.
  \item \textbf{Proverb / saying} (e.g.\ 2017.12, 2020.07, 2020.09): explain the meaning of the saying, illustrate it with an example, and conclude with what you have learned.
\end{enumerate}

\begin{engbox}[title=Universal three-paragraph skeleton]
\textbf{Paragraph 1 (2--3 sentences):} echo the given sentence or introduce the topic; state your thesis.\\
\textbf{Paragraph 2 (5--7 sentences):} two or three supporting points, each with a reason or example, linked by ``To begin with / In addition / What's more''.\\
\textbf{Paragraph 3 (2--3 sentences):} restate your view and end with a call for action or a prediction.\\
Word target: 160--190 words; keep the given opening sentence exactly as it appears.
\end{engbox}

\subsection{Sample Essays}

\subsubsection*{Sample 1: Digital Literacy (2024.06, Set 1 style)}

\begin{engbox}
There is a growing awareness of the importance of digital literacy and skills in today's world. As technology reshapes nearly every aspect of life, the ability to use digital tools wisely has become as basic as reading and writing.

To begin with, digital skills determine how efficiently we learn. Students who can search, evaluate and organize online information are able to turn the Internet into a personal library, while those who cannot may easily get lost in a sea of unreliable content. In addition, digital literacy protects us from risks. Knowing how to manage privacy settings and recognize online fraud helps us avoid losing money or personal data. Finally, such skills are increasingly required in the workplace, where remote meetings, cloud documents and data analysis are routine tasks.

In conclusion, digital literacy is no longer optional but essential. It is high time that we took the initiative to sharpen these skills so that technology serves us rather than controls us.
\end{engbox}

\subsubsection*{Sample 2: Meeting Challenges (2020.12 style)}

\begin{engbox}
Why should students be encouraged to develop the ability to meet challenges? In an age of rapid change, the capacity to face difficulties calmly and creatively matters more than ever.

First, challenges are the best teachers. Each difficulty we overcome leaves us with new experience and stronger confidence, just as muscles grow only under resistance. Second, the ability to meet challenges prepares us for the future. Graduates today may change jobs several times, and only those who can adapt quickly will thrive. Moreover, facing challenges together builds teamwork: when students work on a difficult project, they learn to communicate, compromise and support each other.

Therefore, universities should create more opportunities---such as group projects and internships---for students to practise meeting challenges. What we gain in the process will benefit us throughout our lives.
\end{engbox}

\subsubsection*{Sample 3: A Sense of Responsibility (2019.12 style)}

\begin{engbox}
A sense of community responsibility plays an irreplaceable role in both personal growth and social progress. It means that we care about the common good and are willing to act for it.

On a personal level, taking responsibility gives life meaning. When we volunteer in a community, help an elderly neighbour, or protect the environment around us, we feel needed and valued. On a social level, responsibility is the glue of a harmonious society. If every citizen merely minds his own business, public problems such as waste and pollution will never be solved. Take the ``Clean Plate Campaign'' as an example: millions of small acts of self-discipline have greatly reduced food waste across the country.

In short, community responsibility is a virtue worth practising every day. As long as each of us contributes our part, our society will become a better place to live in.
\end{engbox}

\subsection{Writing Checklist (self-check before handing in)}

\begin{itemize}
  \item Did I copy the given opening sentence \textbf{exactly}, without adding or dropping words?
  \item Is my essay 150--200 words, and does each paragraph focus on one idea?
  \item Have I used at least \textbf{three} linking devices and \textbf{two} complex structures (a clause, a participle phrase, or an emphatic / inverted sentence)?
  \item Have I avoided the most common slips---agreement, tense, singular vs.\ plural, and Chinese-style word order?
  \item Is the handwriting clear? (In the real exam, neat handwriting still earns impression points.)
\end{itemize}
